% ============================================================
% DRAFT — moving evaluation.tex's "Discussion" section (RQ1, RQ2,
% Limitations) into conclusion.tex, per your call: that section
% already answers the research questions and lists what went
% wrong, which is what the Conclusion chapter is for, not something
% that belongs mid-evaluation-chapter.
%
% No cross-references anywhere else in the report point at
% sec:discussion, and none of conclusion.tex's four section labels
% (sec:answers, sec:contributions, sec:limitations, sec:future_work)
% are referenced anywhere either, so this move is safe — nothing
% else needs updating for a broken \ref.
%
% This file is a draft for review, not a finished edit. Nothing here
% has been written into either actual chapter file. Three blocks:
%   M1. DELETE the whole \mysection{Discussion}{...} block from the
%       end of evaluation.tex.
%   M2. ADD the RQ1/RQ2 content to conclusion.tex's (currently empty)
%       "Answers to the Research Questions" section.
%   M3. ADD the (now expanded) Limitations content to conclusion.tex's
%       (currently empty) "Limitations" section.
% I also used the corrected RQ2 wording from
% evaluation_rq2_ctx_depth_sweep_findings.tex (E5) rather than the
% version currently sitting in evaluation.tex, since that's the one
% built on the depth-120 data rather than the superseded depth-40
% sweep — no point moving text into the Conclusion chapter that's
% about to go stale anyway. If you haven't applied that edit yet,
% apply it first (or just use the RQ2 wording below, it's the same).
%
% UPDATE (this revision): you asked for the Limitations block to also
% cover (a) the camera/lighting problem, (b) a stronger statement that
% level~5 wasn't fully tested, (c) the cloud-model comparison's limited
% scope, and (d) the "probe chains" finding. M3 below now has four
% named sub-paragraphs instead of one: the original evidence-base
% paragraph (lightly strengthened on level~5 with exact counts), plus
% three new ones. See the note right before M3 for what "probe chains"
% actually are, since your message asked.
%
% One thing I did NOT do, flagging it rather than just doing it: the
% very last clause of the RQ2 paragraph below ("...which argues for a
% capability and coverage check on the plan before departure as the
% most useful next improvement") reads as a Future Work item sitting
% inside Answers. Conclusion.tex has an empty
% \mysection{Recommendations for Future Work} section right below
% Limitations — if you want, that clause could become the first line
% there instead of staying in Answers. Said as an option, not done
% for you, since trimming the Answers paragraph is your call to make.
%
% SEPARATE QUESTION you asked ("what about this section in the
% evaluation" re: sec:rq2_results intro + Table~\ref{tab:rq2_depth} +
% "Where execution failed" + "Probe chains"): that section stays in
% evaluation.tex untouched. It's raw results (the depth-1-to-12 table
% and what happened in each failing run), not synthesis, so it belongs
% with the methodology and the other results tables, not in the
% Conclusion. The only thing wrong with it is a typo, unrelated to this
% move: line ~241, "every plan visited the expected andmarks in the
% expected order" should read "landmarks". Worth fixing whenever you're
% next in that file, not urgent.
% ============================================================


% ---------- M1. DELETE from chapters/evaluation.tex ----------
% DELETE this entire block (it currently sits at the very end of the
% file, after "\mysection{Reference comparison against cloud models}"):
%
%   \mysection{Discussion}{Discussion}
%   \label{sec:discussion}
%
%   \paragraph{RQ1.} On the physical rover, the local pipeline handled levels 0 to 2
%   without help in all ten runs, and levels 3 to 5 ended in a completed mission in
%   11 of 13 runs (nine directly and two after guidance), with one run halted after guidance at
%   level~3 and one aborted at level~4. Its main weakness at level~4 was dropping the condition
%   in two of seven plans, and its main weakness at level~3 was speech-to-text noise, which twice
%   changed or lost the meaning of the command. Reliability at levels 3 to 5 is therefore supported
%   by a handful of runs and should be reported as a demonstration of capability, not as a rate.
%   The part of the design that earns its place is the guidance loop, which turned several runs
%   that would have ended with ``target not found'' into completed missions.
%
%   \paragraph{RQ2.} In the virtual tests no hallucinated step appeared
%   in chains of up to 12 named landmarks or 8 literal instructions,
%   and chains of up to 36 instructions stayed correct at the configured context.
%   The first failure mode is not hallucination but truncation, at about 34 steps,
%   followed by the silent dropping of steps that cannot be carried out. Neither
%   failure is caught by the pipeline at plan time, which argues for a capability
%   check on the plan before departure as the most useful next improvement.
%
%   \paragraph{Limitations.} The evidence is one room, one rover, and
%   three to seven runs per physical level. Level~5 is not yet adequately
%   covered, the level~4 branches have not been scored against the true scene,
%   and the vision-based distance estimate was not good enough to use.
%   The virtual tests switch odometry drift off and use typed text,
%   so they separate plan quality from perception and speech, and they are
%   not directly comparable to the physical runs. The ultrasonic stop does
%   not protect against low obstacles approached at an angle.
%
% After deleting it, evaluation.tex should end at
% "...no cloud run was made on the physical rover." (the last line of
% \mysection{Reference comparison against cloud models}).


% ---------- M2. ADD to chapters/conclusion.tex, under "Answers to the Research Questions" ----------
% CURRENT:
%   \mysection{Answers to the research questions}{Answers to the Research Questions}
%   \label{sec:answers}
%
%   \mysection{Contributions}{Contributions}
% INSERT the following between those two lines:

\mysection{Answers to the research questions}{Answers to the Research Questions}
\label{sec:answers}

\paragraph{RQ1.} On the physical rover, the local pipeline handled levels 0 to 2
without help in all ten runs, and levels 3 to 5 ended in a completed mission in
11 of 13 runs (nine directly and two after guidance), with one run halted after
guidance at level~3 and one aborted at level~4. Its main weakness at level~4 was
dropping the condition in two of seven plans, and its main weakness at level~3
was speech-to-text noise, which twice changed or lost the meaning of the
command. Reliability at levels 3 to 5 is therefore supported by a handful of
runs and should be reported as a demonstration of capability, not as a rate.
The part of the design that earns its place is the guidance loop, which turned
several runs that would have ended with ``target not found'' into completed
missions.

\paragraph{RQ2.} In the virtual tests no hallucinated step appeared at any
depth tested, including the extended context sweep run out to 120 chained
instructions: across all 270 of those runs, the planner never substitutes a
wrong object, location or action. Compositional depth does have a hard limit,
but the failure mode is truncation, not hallucination: from depth 40 onward,
every run at every context setting tested (1024, 2048 and 8192 tokens) stops
generating steps partway through an otherwise-correct plan, and none recovers
at any depth up to 120. Context window size does not move this onset, but it
does change its severity: the smaller context settings produce a completely
empty plan in roughly 40\% of deep runs, while the largest setting never does
and sustains close to double the plan length on average, which points to a
generation-length ceiling in the planner's own decoding behaviour rather than
the model running out of context-window capacity to track the room and the
command. Neither this truncation nor the silent dropping of unresolvable steps
is caught by the pipeline at plan time, which argues for a capability and
coverage check on the plan before departure as the most useful next
improvement.


% ---------- M3. ADD to chapters/conclusion.tex, under "Limitations" ----------
% What "probe chains" are, since you asked: they're one of the four
% kinds of virtual test chain (Section~\ref{ssec:corpus}), each eight
% steps long, where one step is deliberately impossible — a landmark
% that isn't in the room, or an action outside the capability
% vocabulary ("pick up the lamp"). The point of a probe is to see
% whether the pipeline notices and flags the bad step, or just quietly
% carries on as if it weren't there. Section~\ref{sec:rq2_results}
% ("Probe chains") found it almost always does the latter: in all 9
% cloud runs and 12 of 15 local runs, the planner silently dropped the
% impossible step and returned a shorter, otherwise-valid-looking plan,
% with no record anywhere that a step had to be dropped. That silent-
% dropping behaviour, and the fact that your logging never caught it,
% is now its own Limitations paragraph below ("Silent plan
% modification").
%
% CURRENT:
%   \mysection{Limitations}{Limitations}
%   \label{sec:limitations}
%
%   \mysection{Recommendations for future work}{Recommendations for Future Work}
% INSERT the following between those two lines:

\mysection{Limitations}{Limitations}
\label{sec:limitations}

\paragraph{Evidence base.} The evidence is one room, one rover, and three to
seven runs per physical level. Level~5 was the least settled of the six: of
its ten runs, three completed outright, three only completed after a guidance
turn, and four halted or aborted. That is the most mixed result of any level,
so I do not consider it adequately covered yet. The level~4 branches have not
been scored against the true scene either, and the vision-based distance
estimate was not good enough to use. The virtual tests switch odometry drift
off and use typed text, so they separate plan quality from perception and
speech, and they are not directly comparable to the physical runs. The
ultrasonic stop does not protect against low obstacles approached at an
angle.

\paragraph{Camera and lighting.} The camera, not the planner, was the weak
point of the physical setup. Across levels~3 and~5, every halt with an
operator note attached names poor lighting or low contrast as the reason: a
black bag invisible against a dark floor, an orange box blending into a
cupboard door, a frame thrown out as too blurry because the room was too dim,
and one run that only worked once the operator physically moved the target
into better light. Level~5 shows this best: the plans themselves were judged
good in every run I reviewed, and the failures were the vision model running
out of usable light once the rover got close enough to need a fine judgement.
This is a limitation of the webcam-and-lighting setup I tested with, not of
the planning pipeline. A better-lit room, or a camera that actually works in
low light, would probably fix most of it.

\paragraph{Reference comparison scope.} The cloud-model comparison
(Section~\ref{sec:cloud_comparison}) only covers the depth-1-to-12 virtual
chains. It was never run on the depth-4-to-120 context sweep, and I never ran
a cloud model on the physical rover. So the finding that a small local model
matches a cloud model on correctness and latency only holds for short virtual
chains. It says nothing about the compositional-depth ceiling, or about real
perception and speech.

\paragraph{Silent plan modification.} The probe chains showed that when a
step can't be carried out, the pipeline's default behaviour is to just drop
it from the plan without saying so: in all 9 cloud runs and 12 of 15 local
runs, not one got flagged to the operator. None of those dropped steps was
recorded as unresolved by the pipeline's own logging either
(\texttt{probe\_logged} came back false every time), so right now the system
has no way to tell a silently shortened plan apart from one that actually
finished. That is as much a gap in my logging as it is in the plan itself,
and it is the clearest concrete thing the capability check proposed under RQ2
above should catch.
